\scanpage{53}
\mspage{9}

\uline{NB} Ayant (via un groupoïde équivalent \deleted{\unsure{à}} \unsure{description|de conception} purement combinatoire) ``reconstruit'' ou \deleted{\unsure{remplacé}} réalisé le groupoïde $S\unsure{\mathscr T|\mathbb T}{\unsure{\mathbb D|D}}_{\mathrm{top}}^{\pm}$, et le foncteur
\[
 S\mathscr T\mathbb D_{\mathrm{top}}^{\pm}\longrightarrow(\mathrm{Gr.\ pond.\ rep})
\]

% slanted
\margin{\uline{NB} Cette opération a l'air d'un \uline{affaiblissement} de structure --- mais \unsure{elle|ce} n'est affaiblissement, sur une composante $C$ donnée de $\Sigma^+$, que lorsqu'on ``gomme'' \uline{tous} les repères de $C$ --- sinon c'est au contraire plutôt dans la nature d'un renforcement de structure, --- ou plutôt, du moment qu'on se donne déjà $R'_C=R'\cap C$, dans les catégories isotopiques \unsure{peu importe}, l'opération de gommage et celle \deleted{\unsure{au sens où}} de \unsure{surmarquage|sous-marquage} (cf plus bas) \unsure{correspondantes|corresp.} sont inverses l'une de l'autre, et sont des équivalences de groupoïdes. Il faudrait séparer les deux cas, pour distinguer les \uline{opérations de gommage essentielles} (qui changent \unsure{les groupes|le groupe} des automorphismes) et les autres.}

\deleted{\unsure{cela}} on aura \deleted{\unsure{par}} donc aussi \deleted{une traduction} \inserted{l'équivalent de} \deleted{de} la donnée d'un $(X,\Sigma,R)$, muni \uline{en plus} d'une partie $R'\subset R$ --- \deleted{\unsure{d'où}} i.e. \unsure{on va|on a|on aura} la construction d'un groupoïde \unsure{classifiant|fibrant}
\diagram{\begin{tikzcd}[ampersand replacement=\&]
 S_{*}\mathscr T\mathbb D_{\mathrm{top}}^{\pm}\ \text{(à sous ens. marqué $R'$ de sommets-repères)}\arrow[d]\\
 S_{*}\mathscr T\mathbb D_{\mathrm{top}}^{\pm}
\end{tikzcd}}

et \deleted{\unsure{d'où}} il nous faudra donc \inserted{décrire} un foncteur ``gommage de $R'$''
\diagram{\begin{tikzcd}[ampersand replacement=\&,column sep=large]
 S_{*}\mathscr T\mathbb D_{\mathrm{top}}^{\pm}\bigl(\underset{\text{marqué}}{\overset{\text{sous}}{\deleted{\unsure{R}}R'}}\bigr)
 \arrow[r,"\text{gommage de }R'", "\scriptstyle\rightarrow"{description,pos=.53,inner sep=0pt}]\&S_{*}\mathscr T\mathbb D_{\mathrm{top}}^{\pm}
\end{tikzcd}}

\deleted{Cela s'exprime, sur le} donnant lieu à un diagramme comm. de foncteurs

\diagram{\begin{tikzcd}[ampersand replacement=\&,column sep=small,row sep=small]
S_{*}\mathscr T\mathbb D_{\mathrm{top}}^{\pm}(\ldots R'\ldots)
 \arrow[rr,"\text{gommage de }R'"]\arrow[d,"\text{oubli du }\unsure{\text{marquage}|\text{sous-groupe}|\text{distinguage}}\ R'"] \& \& S_{*}\mathscr T\mathbb D_{\mathrm{top}}^{\pm}\arrow[ddl]\\
S_{*}\mathscr T\mathbb D\deleted{\unsure{R}}_{\mathrm{top}}^{\pm}\arrow[dr] \& \&\\
\&(\mathrm{Gra}\deleted{\unsure{\mathrm{bi}}}\mathrm{pond})\arrow[r] \& \substack{\text{graphes avec}\\\text{pondération}\\\text{par $\gamma$, mais}\\\text{non par $r$,}}
\end{tikzcd}}

alors qu'au niveau
% archive sheet number: 52

